% Figure files are generated with plot_section43_mechanism.py (matplotlib).
% Paste into Section 4.3 after reviewing section43_mechanism.md.
\begin{figure}[t]
  \centering
  \includegraphics[width=\textwidth]{figures/section43_gradient_cosine.pdf}
  \caption{Layerwise gradient alignment under endpoint read noise at trained
  checkpoints. Columns show Conv1--Conv3 and rows show the baseline, balanced,
  and legacy amplification schemes. Each cell is the mean cosine similarity
  between the centered EqProp weight-gradient estimate and a noiseless BPTT
  reference, computed separately for each minibatch and noise realization.
  Both estimators start from the same state after $T$ free iterations; BPTT
  differentiates through the following $K$ iterations. We replay the clean
  seed-0 EqProp best-validation checkpoint on the same 36 validation minibatches
  of 16 examples, with $T=K=4,6,8$ for Conv1--Conv3. Nonzero noise levels use eight
  draws per minibatch, paired across schemes and noise levels, with independent
  Gaussian noise at the positive and negative endpoints. Noise-axis labels
  use a common multiplier of $10^{-5}$. The zero-noise column
  uses one clean estimate per minibatch. All calculations preserve float64
  precision, and weights remain fixed across noise levels. Undefined cosines
  would be shown as dashes. These are validation-set mechanism measurements,
  not test accuracy or averages over training seeds.}
  \label{fig:eqprop-read-noise-gradient-cosine}
\end{figure}

\begin{figure}[t]
  \centering
  \includegraphics[width=\textwidth]{figures/section43_phase_displacement.pdf}
  \caption{Clean phase contrast from the input side to the output. For the same
  checkpoints and validation cohort as Fig.~\ref{fig:eqprop-read-noise-gradient-cosine},
  each curve shows the RMS centered voltage contrast $(v^+-v^-)/2$, pooled over
  all examples and state coordinates in each layer. H1 is the first hidden layer;
  the clamped input is excluded. Voltages are shown in absolute simulator units
  on a shared logarithmic scale. The gray band spans $\sigma/\sqrt{2}$ for the five
  nonzero read-noise levels, the standard deviation of the centered voltage
  contrast when positive and negative endpoint noises are independent. This
  band is a voltage-noise reference, not an exact gradient signal-to-noise ratio
  or a training-failure threshold. Centered contrast isolates the difference
  between the nudged phases; raw free-to-nudged displacement and matched
  zero-nudge drift are retained in the supporting data. Injected nudge strengths
  for baseline/balanced/legacy are $100/30/3$ for Conv1, $100/10/0.03$ for Conv2,
  and $10/3/0.001$ for Conv3, matching the source training configurations.}
  \label{fig:eqprop-layerwise-phase-contrast}
\end{figure}
